\documentclass[10pt,a4paper]{article}
\usepackage{nexusS5}
\setnexusheader{Entrée en Première — Numérique et sciences informatiques — Stage de pré-rentrée}{Évaluation finale — Ahmad BELDI}
\begin{document}
\nexuscover{ÉVALUATION FINALE --- DIAGNOSTIC DE SORTIE}{Durée : 45 minutes}{Ahmad BELDI}{Entrée en Première — Numérique et sciences informatiques}{NSI --- à traiter individuellement}
\par\noindent\begin{tabularx}{\linewidth}{>{\bfseries}p{22mm} X >{\bfseries}p{16mm} X}
\toprule Nom et prénom & \rule{62mm}{0.32pt} & Date & \rule{34mm}{0.32pt} \\ \bottomrule
\end{tabularx}

\begin{goalbox}
\textbf{Ce document est un diagnostic de sortie, pas un classement.} Il mesure ce qui est
désormais installé, ce qui reste fragile, et il sert à écrire le plan des quatre premières
semaines de la rentrée.
\end{goalbox}

\begin{methodbox}
\textbf{Consignes}
\begin{itemize}
\item Durée : \textbf{45 minutes}. Partie A : environ 10 min. Partie B : environ 20 min. Partie C : environ 15 min.
\item Travailler seul, sans document et sans les cartes de rappel de la séance.
\item Écrire la relation ou la propriété utilisée avant le calcul.
\item Une question peut être laissée de côté puis reprise ; il vaut mieux la laisser vide que d'écrire au hasard.
\item Trois lignes « Certitude » seulement : \conf\ , sur la même échelle qu'en début de stage.
\end{itemize}
\end{methodbox}

\newpage
\phasehead{Partie A --- Automatismes et connaissances essentielles}{environ 10 min}{Questions courtes. Écrire le résultat, sans rédaction longue.}
\begin{enumerate}
\needspace{22mm}
\item[\textbf{A1.}] Pour \code{L = [7, 1, 8, 4, 2]}, donner \code{L[2]}, \code{L[-1]}, \code{len(L)} et le dernier indice valide.
\worklines{2}
\par\vspace{1.2mm}
\needspace{22mm}
\item[\textbf{A2.}] Écrire la négation de \code{x <= 8}, puis celle de \code{(x > 2) or (x == 0)}.
\worklines{2}
\par\vspace{1.2mm}
\needspace{22mm}
\item[\textbf{A3.}] Pour \code{for i in range(3, 12, 4):}, donner les valeurs prises par \code{i} et le nombre de tours.
\worklines{2}
\par\vspace{1.2mm}
\needspace{22mm}
\item[\textbf{A4.}] On exécute \code{a = 6}, puis \code{b = a}, puis \code{a = a + 4}. Donner \code{a} et \code{b}, puis expliquer en une phrase.
\worklines{2}
\par\vspace{1.2mm}
\needspace{22mm}
\item[\textbf{A5.}] On définit \code{def affiche(n): print(n * 2)}. Que vaut \code{r} après \code{r = affiche(3)} ? Expliquer en une phrase.
\worklines{2}
\par\vspace{1.2mm}
\needspace{22mm}
\item[\textbf{A6.}] Donner le nombre de tours de \code{for i in range(6)}, puis les valeurs prises par \code{i} dans \code{for i in range(2, 9)}.
\worklines{2}
\par\vspace{1.2mm}
\end{enumerate}
\certbox{Certitude sur l'ensemble de la partie A :}
\needspace{68mm}\par\vspace{3mm}
\phasehead{Partie B --- Application et raisonnement}{environ 20 min}{Écrire la relation, la propriété ou la règle avant le calcul, puis contrôler.}
\begin{enumerate}
\needspace{22mm}
\item[\textbf{B1.}] On exécute le programme suivant.
\begin{lstlisting}
total = 0
n = 0
for v in [4, 9, 2, 7]:
    total = total + v
    n = n + 1
\end{lstlisting}
\begin{enumerate}
\item Compléter le tableau d'état après chaque tour de boucle.
\par\noindent\begin{tabularx}{\linewidth}{|>{\bfseries}p{16mm}|X|X|X|X|}
\hline
Après le tour & 1 & 2 & 3 & 4 \\ \hline
\code{total} & & & & \\ \hline
\code{n} & & & & \\ \hline
\end{tabularx}
\item Donner les valeurs finales, puis dire en une phrase ce que représente chacune des deux variables.
\end{enumerate}
\worklines{4}
\par\vspace{1.2mm}
\needspace{22mm}
\item[\textbf{B2.}] On dispose de la fonction suivante.
\begin{lstlisting}
def moyenne(valeurs):
    total = 0
    for v in valeurs:
        total = total + v
    print(total / len(valeurs))
\end{lstlisting}
\begin{enumerate}
\item On écrit \code{m = moyenne([4, 6])}. Que vaut \code{m} ? Justifier.
\item Corriger la fonction pour qu'elle \emph{renvoie} la moyenne.
\end{enumerate}
\worklines{5}
\par\vspace{1.2mm}
\needspace{22mm}
\item[\textbf{B3.}] On dispose de l'extrait de fichier suivant.
\begin{lstlisting}[language={}]
id,date,temperature,humidite,statut
C07,2026-08-24,26.5,58,OK
C08,2026-08-24,32.1,49,ALERTE
\end{lstlisting}
\begin{enumerate}
\item Quelles colonnes doivent être converties avant tout calcul, et en quel type ?
\item Expliquer pourquoi \code{"26.5" + "1"} ne donne pas $27{,}5$.
\item Donner un exemple de \emph{métadonnée} associable à ce fichier, et dire en quoi elle
diffère d'une donnée du tableau.
\end{enumerate}
\worklines{5}
\par\vspace{1.2mm}
\needspace{22mm}
\item[\textbf{B4.}] On pose \code{L = [10, 4, 7]}.
\begin{enumerate}
\item Donner \code{L[0]}, \code{len(L)} et le dernier indice valide.
\item On exécute ensuite \code{M = L} puis \code{M.append(9)}. Que contient \code{L} ? Expliquer précisément.
\end{enumerate}
\worklines{6}
\par\vspace{1.2mm}
\end{enumerate}
\certbox{Certitude sur la question B2 :}
\needspace{68mm}\par\vspace{3mm}
\phasehead{Partie C --- Transfert et synthèse}{environ 15 min}{Reconnaître la notion utile, mobiliser plusieurs acquis, conclure par une phrase.}
\begin{enumerate}
\needspace{22mm}
\item[\textbf{C1.}] \textbf{Rechercher un capteur.} On dispose d'une liste de dictionnaires
\code{mesures}, chacun possédant au moins la clé \code{"id"}.
\begin{lstlisting}
def chercher(mesures, identifiant):
    for i in range(len(mesures)):
        if ____________________________:
            return ____
    return -1
\end{lstlisting}
\begin{enumerate}
\item Compléter les deux trous.
\item On pose \code{mesures = [{"id": "C01"}, {"id": "C02"}, {"id": "C03"}]}. Donner le résultat de
\code{chercher(mesures, "C03")} et de \code{chercher(mesures, "C09")}, puis écrire un test portant sur
une liste vide, avec son résultat attendu.
\item La liste est désormais triée par identifiant. En une phrase, dire ce que permet alors la recherche
dichotomique et à quelle condition.
\end{enumerate}
\worklines{8}
\par\vspace{1.2mm}
\needspace{22mm}
\item[\textbf{C2.}] Compléter la fonction qui renvoie l'indice de la première valeur strictement négative
d'une liste, ou $-1$ s'il n'y en a aucune.
\begin{lstlisting}
def premier_negatif(L):
    for i in range(len(L)):
        if ______________:
            return ____
    return ____
\end{lstlisting}
Donner ensuite le résultat de \code{premier_negatif([4, 7, -2, -5])}.
\worklines{6}
\par\vspace{1.2mm}
\end{enumerate}
\certbox{Certitude sur la question C1 :}
\brouillon{\dimexpr\textheight/3\relax}
\newpage
\sectiontitle{Après l'évaluation --- à remplir en deux minutes}
\begin{methodbox}Ces trois lignes ne sont pas notées. Elles servent à construire le plan des quatre premières semaines de la rentrée.\end{methodbox}
\par\noindent\begin{tabularx}{\linewidth}{>{\bfseries}p{62mm} X}
\toprule
La question que j'ai trouvée la plus sûre & \rule{85mm}{0.32pt} \\
La question sur laquelle j'ai le plus hésité & \rule{85mm}{0.32pt} \\
Une notion que je veux revoir dès la première semaine & \rule{85mm}{0.32pt} \\
\bottomrule\end{tabularx}
\vfill
\begin{graybox}\small Ce document est un diagnostic de sortie. Il sera analysé compétence par compétence, puis comparé au positionnement passé en début de stage lorsque cette comparaison est légitime.\end{graybox}
\end{document}
